% Options for packages loaded elsewhere
\PassOptionsToPackage{unicode}{hyperref}
\PassOptionsToPackage{hyphens}{url}
\documentclass[
]{article}
\usepackage{xcolor}
\usepackage{amsmath,amssymb}
\setcounter{secnumdepth}{-\maxdimen} % remove section numbering
\usepackage{iftex}
\ifPDFTeX
  \usepackage[T1]{fontenc}
  \usepackage[utf8]{inputenc}
  \usepackage{textcomp} % provide euro and other symbols
\else % if luatex or xetex
  \usepackage{unicode-math} % this also loads fontspec
  \defaultfontfeatures{Scale=MatchLowercase}
  \defaultfontfeatures[\rmfamily]{Ligatures=TeX,Scale=1}
\fi
\usepackage{lmodern}
\ifPDFTeX\else
  % xetex/luatex font selection
\fi
% Use upquote if available, for straight quotes in verbatim environments
\IfFileExists{upquote.sty}{\usepackage{upquote}}{}
\IfFileExists{microtype.sty}{% use microtype if available
  \usepackage[]{microtype}
  \UseMicrotypeSet[protrusion]{basicmath} % disable protrusion for tt fonts
}{}
\makeatletter
\@ifundefined{KOMAClassName}{% if non-KOMA class
  \IfFileExists{parskip.sty}{%
    \usepackage{parskip}
  }{% else
    \setlength{\parindent}{0pt}
    \setlength{\parskip}{6pt plus 2pt minus 1pt}}
}{% if KOMA class
  \KOMAoptions{parskip=half}}
\makeatother
\usepackage{color}
\usepackage{fancyvrb}
\newcommand{\VerbBar}{|}
\newcommand{\VERB}{\Verb[commandchars=\\\{\}]}
\DefineVerbatimEnvironment{Highlighting}{Verbatim}{commandchars=\\\{\}}
% Add ',fontsize=\small' for more characters per line
\newenvironment{Shaded}{}{}
\newcommand{\AlertTok}[1]{\textcolor[rgb]{1.00,0.00,0.00}{\textbf{#1}}}
\newcommand{\AnnotationTok}[1]{\textcolor[rgb]{0.38,0.63,0.69}{\textbf{\textit{#1}}}}
\newcommand{\AttributeTok}[1]{\textcolor[rgb]{0.49,0.56,0.16}{#1}}
\newcommand{\BaseNTok}[1]{\textcolor[rgb]{0.25,0.63,0.44}{#1}}
\newcommand{\BuiltInTok}[1]{\textcolor[rgb]{0.00,0.50,0.00}{#1}}
\newcommand{\CharTok}[1]{\textcolor[rgb]{0.25,0.44,0.63}{#1}}
\newcommand{\CommentTok}[1]{\textcolor[rgb]{0.38,0.63,0.69}{\textit{#1}}}
\newcommand{\CommentVarTok}[1]{\textcolor[rgb]{0.38,0.63,0.69}{\textbf{\textit{#1}}}}
\newcommand{\ConstantTok}[1]{\textcolor[rgb]{0.53,0.00,0.00}{#1}}
\newcommand{\ControlFlowTok}[1]{\textcolor[rgb]{0.00,0.44,0.13}{\textbf{#1}}}
\newcommand{\DataTypeTok}[1]{\textcolor[rgb]{0.56,0.13,0.00}{#1}}
\newcommand{\DecValTok}[1]{\textcolor[rgb]{0.25,0.63,0.44}{#1}}
\newcommand{\DocumentationTok}[1]{\textcolor[rgb]{0.73,0.13,0.13}{\textit{#1}}}
\newcommand{\ErrorTok}[1]{\textcolor[rgb]{1.00,0.00,0.00}{\textbf{#1}}}
\newcommand{\ExtensionTok}[1]{#1}
\newcommand{\FloatTok}[1]{\textcolor[rgb]{0.25,0.63,0.44}{#1}}
\newcommand{\FunctionTok}[1]{\textcolor[rgb]{0.02,0.16,0.49}{#1}}
\newcommand{\ImportTok}[1]{\textcolor[rgb]{0.00,0.50,0.00}{\textbf{#1}}}
\newcommand{\InformationTok}[1]{\textcolor[rgb]{0.38,0.63,0.69}{\textbf{\textit{#1}}}}
\newcommand{\KeywordTok}[1]{\textcolor[rgb]{0.00,0.44,0.13}{\textbf{#1}}}
\newcommand{\NormalTok}[1]{#1}
\newcommand{\OperatorTok}[1]{\textcolor[rgb]{0.40,0.40,0.40}{#1}}
\newcommand{\OtherTok}[1]{\textcolor[rgb]{0.00,0.44,0.13}{#1}}
\newcommand{\PreprocessorTok}[1]{\textcolor[rgb]{0.74,0.48,0.00}{#1}}
\newcommand{\RegionMarkerTok}[1]{#1}
\newcommand{\SpecialCharTok}[1]{\textcolor[rgb]{0.25,0.44,0.63}{#1}}
\newcommand{\SpecialStringTok}[1]{\textcolor[rgb]{0.73,0.40,0.53}{#1}}
\newcommand{\StringTok}[1]{\textcolor[rgb]{0.25,0.44,0.63}{#1}}
\newcommand{\VariableTok}[1]{\textcolor[rgb]{0.10,0.09,0.49}{#1}}
\newcommand{\VerbatimStringTok}[1]{\textcolor[rgb]{0.25,0.44,0.63}{#1}}
\newcommand{\WarningTok}[1]{\textcolor[rgb]{0.38,0.63,0.69}{\textbf{\textit{#1}}}}
\usepackage{longtable,booktabs,array}
\usepackage{caption}
\captionsetup[table]{skip=6pt}
\newcounter{none} % for unnumbered tables
\usepackage{calc} % for calculating minipage widths
% Correct order of tables after \paragraph or \subparagraph
\usepackage{etoolbox}
\makeatletter
\patchcmd\longtable{\par}{\if@noskipsec\mbox{}\fi\par}{}{}
\makeatother
\setlength{\emergencystretch}{3em} % prevent overfull lines
\providecommand{\tightlist}{%
  \setlength{\itemsep}{0pt}\setlength{\parskip}{0pt}}
\usepackage{bookmark}
\IfFileExists{xurl.sty}{\usepackage{xurl}}{} % add URL line breaks if available
\urlstyle{same}
% fallback for those not using the hyperref driver hyperxmp:
\makeatletter
\@ifundefined{xmpquote}{\newcommand{\xmpquote}[1]{#1}}{}
\makeatother
\hypersetup{
  hidelinks,
  pdfcreator={LaTeX via pandoc}}

\author{}
\date{}

\begin{document}

{
\setcounter{tocdepth}{2}
\tableofcontents
}
\section{Block Programming on WebCoder --- Blockly AND
Scratch}\label{block-programming-on-webcoder--blockly-and-scratch}

\begin{quote}
Owner directive (2026-10-30): the platform ships \textbf{both} block
engines --- Blockly \emph{and} Scratch --- not one or the other. Blockly
is the structured bridge to typed code; Scratch is the familiarity layer
(the Scratch cat, the stage, say/answer) for the youngest learners. This
document is the complete guide to both: architecture, judge integration,
authoring, learning, teaching, testing, security, troubleshooting and
worked examples. It is written so a new contributor can extend the
system without reading the code first.
\end{quote}

\begin{center}\rule{0.5\linewidth}{0.5pt}\end{center}

\subsection{1. What shipped, in one
paragraph}\label{1-what-shipped-in-one-paragraph}

A problem can be a \textbf{block problem}. The creator picks a block
\emph{engine} per problem:

\begin{itemize}
\tightlist
\item
  \textbf{Blockly}
  (\texttt{block\_config.engine:\ \textquotesingle{}blockly\textquotesingle{}},
  the default) --- a Blockly editor with a WebCoder-curated toolbox
  (input/control/logic/math/text/
  \textbf{lists}/\textbf{variables}/\textbf{functions}). Learners snap
  blocks; Blockly\textquotesingle s official generators produce
  \textbf{real Python 3 or JavaScript (Node 20)}; that generated source
  is submitted and judged through the \textbf{normal pipeline} ---
  identical limits, identical pipeline, identical verdicts.
\item
  \textbf{Scratch}
  (\texttt{block\_config.engine:\ \textquotesingle{}scratch\textquotesingle{}})
  --- a real Scratch editing experience on site: \textbf{scratch-vm +
  scratch-blocks + scratch-render} (the same libraries
  Scratch\textquotesingle s own GUI uses), with the Scratch cat, a live
  stage, ask/answer input and say output. Learners edit a project; the
  \textbf{project JSON (sb3)} is submitted as the pseudo-language
  \texttt{scratch}, and the judge \textbf{transpiles it to Python}
  (\texttt{backend/problems/scratch\_transpiler.py}) and judges that
  through the normal pipeline under \textbf{python3\textquotesingle s
  language-fair limits}.
\end{itemize}

Both engines are on the free playground at \textbf{\texttt{/blocks}}
(engine switcher at the top). Both localize fully (RO/EN block labels +
UI chrome via the ten-locale i18n catalogs).

\begin{center}\rule{0.5\linewidth}{0.5pt}\end{center}

\subsection{2. Architecture}\label{2-architecture}

\subsubsection{2.1 The two editors
(frontend)}\label{21-the-two-editors-frontend}

\begin{verbatim}
src/components/blocks/
├── BlockEditor.tsx      Blockly layer (PR #252, extended in this wave)
├── wcBlocks.ts          the two site blocks (read number / read line) +
│                        their Python/JS generator snippets
├── complexity.ts        complexity-hint analyzer (pure, tested)
└── ScratchEditor.tsx    Scratch layer (scratch-vm/blocks/render/storage)
\end{verbatim}

\textbf{BlockEditor} lazy-imports \texttt{blockly}
(\texttt{next/dynamic}, \texttt{ssr:\ false} --- SSR never loads it),
registers the two site blocks, applies the official RO or EN message
pack plus the site-block messages, injects a workspace with the
problem\textquotesingle s toolbox (or the default beginner toolbox),
previews the generated code, and \textbf{analyzes the generated code for
complexity hints} (§7). Generation goes through
Blockly\textquotesingle s own \texttt{pythonGenerator} /
\texttt{javascriptGenerator} --- we do not maintain a generator.

\textbf{ScratchEditor} lazy-imports the Scratch stack:

\begin{itemize}
\tightlist
\item
  \texttt{scratch-vm} --- the VM (block representation, threads, green
  flag);
\item
  \texttt{scratch-blocks} --- the Scratch block editor (vertical blocks,
  Scratch look-and-feel). Imported via its self-contained dist bundle
  (\texttt{scratch-blocks/dist/vertical.js}): the
  package\textquotesingle s \texttt{browser} field points at a shim
  built with webpack inline-loader syntax that Turbopack cannot resolve
  --- the dist bundle is the same module, loader-free (§14 covers
  upgrade notes);
\item
  \texttt{scratch-render} --- the WebGL stage;
\item
  \texttt{scratch-storage} --- asset loading, with the \textbf{default
  assets (the cat costume + backdrop) bundled locally} under
  \texttt{public/scratch-assets/} so the starter project loads with zero
  external requests;
\item
  \texttt{scratch-l10n} --- official RO block labels
  (\texttt{scratch-l10n/editor/blocks/ro.json} applied onto
  \texttt{Blockly.Msg}, exactly how scratch-gui localizes).
\end{itemize}

Wiring (same as scratch-gui, minimized):

\begin{enumerate}
\def\labelenumi{\arabic{enumi}.}
\tightlist
\item
  \texttt{storage} created; bundled assets pre-registered in
  \texttt{storage.builtinHelper} under their md5 ids;
\item
  \texttt{vm.attachStorage(storage)};
  \texttt{vm.attachRenderer(new\ ScratchRender(canvas))} (renderer init
  is fault-tolerant: no WebGL ⇒ blank stage, but editing and export keep
  working --- grading never needs the renderer);
\item
  \texttt{vm.start()};
  \texttt{workspace.addChangeListener(vm.blockListener)} --- the
  workspace is the editor, the VM holds the truth;
\item
  \texttt{vm.on(\textquotesingle{}workspaceUpdate\textquotesingle{},\ (\{xml\})\ =\textgreater{}\ Blockly.Xml.clearWorkspaceAndLoadFromXml(...))}
  --- how a loaded project paints its blocks into the workspace;
\item
  \texttt{vm.loadProject(JSON.parse(starter))} then
  \texttt{vm.emitWorkspaceUpdate()};
\item
  \textbf{export}:
  \texttt{vm.on(\textquotesingle{}PROJECT\_CHANGED\textquotesingle{},\ ()\ =\textgreater{}\ onProjectChange(vm.toJSON()))}
  --- \texttt{vm.toJSON()} is scratch-vm\textquotesingle s official sb3
  serializer;
\item
  \textbf{run}: the green-flag button calls \texttt{vm.greenFlag()} (the
  local stage); \texttt{say}/\texttt{think} blocks also stream into a
  \textbf{console strip} under the stage
  (\texttt{vm.runtime.on(\textquotesingle{}SAY\textquotesingle{})}),
  mirroring what the judge will compare.
\end{enumerate}

\textbf{Scratch block media} (block icons) are served from
\texttt{public/scratch-blocks-media/} (copied from
\texttt{node\_modules/scratch-blocks/media/}; re-copy when upgrading
scratch-blocks).

\subsubsection{2.2 The judge integration
(backend)}\label{22-the-judge-integration-backend}

\begin{verbatim}
problem.block_config                     (JSON column, problems/models.py)
   engine: 'blockly' → generator + toolbox + starter (Blockly XML)
   engine: 'scratch' → scratch_project (sb3 JSON) + scratch_blocks (opcodes)

Submission.language
   registry ids (python3, cpp17, …)  → judged directly
   'scratch' (pseudo-language)       → transpiled then judged as python3
\end{verbatim}

\textbf{Blockly submissions} carry the \emph{generated source};
\texttt{language} is the generator (\texttt{python3} /
\texttt{nodejs20}). Nothing special exists in the judge for them --- by
design.

\textbf{Scratch submissions} carry the \emph{sb3 project JSON};
\texttt{language} is \texttt{scratch}. Three enforcement points (the
same defense-in-depth pattern as the wave-C library policy and wave-E
template integrity):

\begin{enumerate}
\def\labelenumi{\arabic{enumi}.}
\tightlist
\item
  \textbf{Submit-time validation} (\texttt{submissions/serializers.py}):
  the problem must be
  \texttt{engine:\ \textquotesingle{}scratch\textquotesingle{}}; the
  project must parse, use only mapped opcodes \textbf{and} only the
  problem\textquotesingle s curated allowlist
  (\texttt{validate\_project}). Violations → HTTP 400 naming the block.
\item
  \textbf{Judge-time re-validation + transpile}
  (\texttt{submissions/tasks.py}): the exact submitted bytes are
  re-validated and transpiled (\texttt{transpile(submission.code)}); a
  failure lands as \textbf{COMPILE\_ERROR} with the reason (a client
  that skipped the API still cannot judge).
\item
  \textbf{Execution}: the transpiled Python is compiled/run under the
  registry\textquotesingle s \texttt{python3} spec; effective limits
  resolve through
  \texttt{judged\_language(\textquotesingle{}scratch\textquotesingle{})\ ==\ \textquotesingle{}python3\textquotesingle{}}
  --- i.e. \textbf{python3\textquotesingle s J5 multipliers (3× time,
  1.5× memory)}, with
  \texttt{Problem.limit\_overrides{[}\textquotesingle{}scratch\textquotesingle{}{]}}
  still able to override either axis. See
  \texttt{submissions/judge\_utils/languages.py}
  (\texttt{SCRATCH\_LANGUAGE\_ID}, \texttt{SCRATCH\_JUDGED\_AS},
  \texttt{judged\_language}).
\end{enumerate}

\subsubsection{2.3 Which judging path shipped and why (decision
record)}\label{23-which-judging-path-shipped-and-why-decision-record}

The owner requirement allowed two Scratch judging paths: headless
scratch-vm execution inside the sandboxed judge, \textbf{or} transpiling
the project to Python via a defined mapping, with the choice documented.
\textbf{The transpiling path shipped.} Reasons:

\begin{enumerate}
\def\labelenumi{\arabic{enumi}.}
\tightlist
\item
  \textbf{Determinism.} scratch-vm is a 30-frames-per-second
  soft-realtime VM designed around a browser stage; its thread
  scheduling is frame-quantized. A judge must measure \emph{CPU time
  against a limit} --- frame-stepped execution makes TLE semantics fuzzy
  and verdicts machine-dependent. The transpiler\textquotesingle s
  output is ordinary Python under the ordinary, deterministic judge.
\item
  \textbf{Language-fair limits compose.} The whole point of the
  wave\textquotesingle s item A is per-language normalization. A
  transpiled Scratch program \emph{is} a CPython program, so it inherits
  python3\textquotesingle s multipliers for free. A headless scratch-vm
  would have needed its own calibration with no external judge practice
  to cite.
\item
  \textbf{Sandbox posture.} The judge-runner container is
  network-isolated with a read-only rootfs; running scratch-vm there
  would mean baking a \textasciitilde250 MB-unpacked node dependency
  tree (scratch-vm pulls scratch-render, scratch-storage, scratch-audio,
  \ldots) into the hardened image --- a large supply-chain and
  image-size cost for a less deterministic outcome. The transpiler is
  \textasciitilde600 lines of pure Python inside the existing worker.
\item
  \textbf{The learner-visible behavior is preserved.}
  \texttt{ask}/\texttt{answer} map to stdin lines;
  \texttt{say}/\texttt{think} map to stdout lines with
  Scratch\textquotesingle s exact number rendering (§5.4). What the
  learner watched on the stage is what the judge compares.
\end{enumerate}

The headless path remains a documented future option: if Scratch
problems ever need stage-dependent semantics (sprite positions, pen,
collisions), a scratch-vm lane in the judge image is the escape hatch
--- the submission format (sb3 JSON) would not change.

\begin{center}\rule{0.5\linewidth}{0.5pt}\end{center}

\subsection{3. Authoring guide
(creators)}\label{3-authoring-guide-creators}

\subsubsection{3.1 Creating a Blockly
problem}\label{31-creating-a-blockly-problem}

\begin{enumerate}
\def\labelenumi{\arabic{enumi}.}
\item
  Create a problem as usual; under \textbf{Block programming} check the
  box.
\item
  Engine \textbf{Blockly}; pick the generator (\textbf{Python 3} or
  \textbf{JavaScript}).
\item
  Optionally paste a \textbf{Toolbox XML} --- the curated block set.
  Omit it to get the default beginner toolbox (logic, loops, math, text,
  input, \textbf{lists}, variables, \textbf{functions}). Copy a category
  subset to curate, e.g. a variables+math-only problem:

\begin{Shaded}
\begin{Highlighting}[]
\NormalTok{\textless{}}\KeywordTok{xml} \OtherTok{xmlns=}\StringTok{"https://developers.google.com/blockly/xml"}\NormalTok{\textgreater{}}
\NormalTok{  \textless{}}\KeywordTok{category} \OtherTok{name=}\StringTok{"Input"} \OtherTok{colour=}\StringTok{"\%\{BKY\_TEXTS\_HUE\}"}\NormalTok{\textgreater{}}
\NormalTok{    \textless{}}\KeywordTok{block} \OtherTok{type=}\StringTok{"wc\_input\_number"}\NormalTok{\textgreater{}\textless{}/}\KeywordTok{block}\NormalTok{\textgreater{}}
\NormalTok{  \textless{}/}\KeywordTok{category}\NormalTok{\textgreater{}}
\NormalTok{  \textless{}}\KeywordTok{category} \OtherTok{name=}\StringTok{"Math"} \OtherTok{colour=}\StringTok{"\%\{BKY\_MATH\_HUE\}"}\NormalTok{\textgreater{}}
\NormalTok{    \textless{}}\KeywordTok{block} \OtherTok{type=}\StringTok{"math\_number"}\NormalTok{\textgreater{}\textless{}/}\KeywordTok{block}\NormalTok{\textgreater{}}
\NormalTok{    \textless{}}\KeywordTok{block} \OtherTok{type=}\StringTok{"math\_arithmetic"}\NormalTok{\textgreater{}\textless{}/}\KeywordTok{block}\NormalTok{\textgreater{}}
\NormalTok{    \textless{}}\KeywordTok{block} \OtherTok{type=}\StringTok{"math\_modulo"}\NormalTok{\textgreater{}\textless{}/}\KeywordTok{block}\NormalTok{\textgreater{}}
\NormalTok{  \textless{}/}\KeywordTok{category}\NormalTok{\textgreater{}}
\NormalTok{  \textless{}}\KeywordTok{category} \OtherTok{name=}\StringTok{"Variables"} \OtherTok{colour=}\StringTok{"\%\{BKY\_VARIABLES\_HUE\}"} \OtherTok{custom=}\StringTok{"VARIABLE"}\NormalTok{\textgreater{}\textless{}/}\KeywordTok{category}\NormalTok{\textgreater{}}
\NormalTok{\textless{}/}\KeywordTok{xml}\NormalTok{\textgreater{}}
\end{Highlighting}
\end{Shaded}

  Site blocks: \texttt{wc\_input\_number} (read an integer) and
  \texttt{wc\_input\_line} (read a text line) --- generated Python
  \texttt{int(input())} / \texttt{input()}; JS \texttt{wcReadNumber()} /
  \texttt{wcReadLine()} with a one-shot stdin-reader preamble the editor
  prepends automatically.
\item
  Optionally paste \textbf{Starter workspace XML} (pre-placed blocks).
\item
  Set \texttt{allowed\_languages} to the generator(s) you accept. The
  generated code is judged exactly like typed code --- including the
  language-fair limits shown on the problem page.
\end{enumerate}

\subsubsection{3.2 Creating a Scratch
problem}\label{32-creating-a-scratch-problem}

\begin{enumerate}
\def\labelenumi{\arabic{enumi}.}
\item
  Create a problem; check \textbf{Block programming}; engine
  \textbf{Scratch}.
\item
  \textbf{Allowed languages}: tick \textbf{Scratch (blocks)} (the
  pseudo-language; its effective limits --- python3\textquotesingle s
  multipliers --- are shown on the checkbox like every other language).
\item
  Optionally paste a \textbf{starter Scratch project} (sb3
  \texttt{project.json}). Default: the Scratch cat with
  \texttt{when\ flag\ clicked\ →\ say\ "Hello!"}. \textbf{Declare the
  variables you want learners to use in the starter} --- v1 offers the
  palette\textquotesingle s variable blocks for pre-declared variables
  (there is no "make a variable" modal in the curated editor; this is
  the curated-set philosophy, not an oversight --- see §9 roadmap).
\item
  Optionally list \textbf{allowed Scratch blocks} (comma-separated
  opcodes). The universe is the judged-mappable set
  (\texttt{KNOWN\_OPCODES} in
  \texttt{backend/problems/scratch\_transpiler.py}):

  {\def\LTcaptype{none} % do not increment counter
  \begin{longtable}[]{@{}ll@{}}
  \toprule\noalign{}
  Category & Opcodes \\
  \midrule\noalign{}
  \endhead
  \bottomrule\noalign{}
  \endlastfoot
  Events & \texttt{event\_whenflagclicked} \\
  Control & \texttt{control\_repeat}, \texttt{control\_repeat\_until},
  \texttt{control\_whileUntil}, \texttt{control\_forever},
  \texttt{control\_if}, \texttt{control\_if\_else},
  \texttt{control\_stop} \\
  Looks & \texttt{looks\_say}, \texttt{looks\_sayforsecs},
  \texttt{looks\_think} \\
  Sensing & \texttt{sensing\_askandwait}, \texttt{sensing\_answer} \\
  Operators & \texttt{operator\_add/subtract/multiply/divide/mod},
  \texttt{operator\_lt/gt/equals}, \texttt{operator\_and/or/not},
  \texttt{operator\_mathop}
  (abs/floor/ceiling/sqrt/sin/cos/tan/round) \\
  Variables & \texttt{data\_variable}, \texttt{data\_setvariableto},
  \texttt{data\_changevariableby} \\
  \end{longtable}
  }

  Leave it empty for the full curated set. \texttt{control\_wait} (no
  judged semantics) and \texttt{operator\_random} (nondeterministic) are
  never offered.
\item
  \textbf{Testcases} are ordinary stdin/stdout pairs: each \texttt{ask}
  consumes one input line; each \texttt{say}/\texttt{think} emits one
  output line.
\end{enumerate}

\subsubsection{3.3 Scratch project packaging (how to produce the starter
JSON)}\label{33-scratch-project-packaging-how-to-produce-the-starter-json}

The starter is the sb3 \textbf{project.json} document (NOT the zipped
.sb3). Ways to produce one:

\begin{itemize}
\tightlist
\item
  \textbf{From this site}: open \texttt{/blocks}, switch the engine to
  \textbf{Scratch}, build the starter (add variables first --- they
  appear in the palette), then take the exported JSON from your
  browser\textquotesingle s network tab
  (\texttt{POST\ /api/v1/submissions/submit/} payload when submitting to
  any scratch problem) --- or ask an admin to fetch it from the
  submission.
\item
  \textbf{From the Scratch editor} (scratch.mit.edu): build the project,
  File → Save to your computer (\texttt{.sb3}), then unzip and use
  \texttt{project.json}. Costumes other than the bundled cat/backdrop
  will load in the learner\textquotesingle s editor \textbf{from the
  official asset CDN} (needs internet on the learner side); the judge
  itself never fetches assets (blocks only).
\end{itemize}

Assets referenced by md5 that exist under
\texttt{public/scratch-assets/} load locally with no network. The
bundled pair: \texttt{bcf454acf82e4504149f7ffe07081dbc.svg} (the cat),
\texttt{cd21514d0531fdffb22204e0ec5ed84a.svg} (the default backdrop).

\subsubsection{3.4 What the creator UI
enforces}\label{34-what-the-creator-ui-enforces}

\texttt{block\_config} is validated at authoring time
(\texttt{problems/serializers.py}): engine ∈ \{blockly, scratch\};
Blockly keys (\texttt{generator}/\texttt{toolbox}/\texttt{starter}) vs
Scratch keys (\texttt{scratch\_project}/\texttt{scratch\_blocks}) must
not mix; \texttt{scratch\_blocks} entries must be judged-mappable
opcodes; sizes are capped (XML 100 KB, project JSON 500 KB).

\begin{center}\rule{0.5\linewidth}{0.5pt}\end{center}

\subsection{4. Learner guide}\label{4-learner-guide}

\begin{itemize}
\tightlist
\item
  \textbf{Blockly}: pick blocks from the categories, connect them, press
  \emph{Show generated code} to see the real Python/JavaScript you built
  --- then submit. It is judged like any submission; the problem page
  shows the effective time/memory for your language. Watch the ⚠ hint
  chips: they appear when the generated code has nested loops (an n²
  surprise) or a loop that never stops.
\item
  \textbf{Scratch}: drag blocks from the palette (the Scratch cat is
  your sprite). Use \texttt{ask\ and\ wait} to read input, \texttt{say}
  to print, variables and operators to compute. Press the \textbf{green
  flag} to run on the stage --- \texttt{say} lines also appear in the
  console strip, which is exactly what the judge compares. Submit sends
  your project; it is translated to Python behind the curtain and judged
  with Python\textquotesingle s time limits.
\end{itemize}

Both playgrounds are free at \texttt{/blocks} --- no submission spent.

\begin{center}\rule{0.5\linewidth}{0.5pt}\end{center}

\subsection{5. The Scratch→Python mapping
(precise)}\label{5-the-scratchpython-mapping-precise}

Implemented in \texttt{backend/problems/scratch\_transpiler.py}
(\textasciitilde100\% line coverage in
\texttt{backend/problems/test\_scratch\_transpiler.py}, including
executing the generated code). Everything outside the mapping raises a
\texttt{TranspileError} naming the opcode --- never a silently-wrong
translation.

\subsubsection{5.1 Program shape}\label{51-program-shape}

\begin{itemize}
\tightlist
\item
  Scripts start at \texttt{when\ green\ flag\ clicked} hats. Multiple
  flag scripts run \textbf{sequentially in stage order} (top-to-bottom,
  left-to-right --- Scratch\textquotesingle s own ordering). Documented
  deviation: Scratch runs them as parallel threads; the console problems
  this layer targets are ask→compute→say scripts that are independent by
  construction.
\item
  \texttt{stop\ all} → \texttt{sys.exit(0)};
  \texttt{stop\ this\ script}/\texttt{other\ scripts} → \texttt{pass}
  (sequential execution makes them equivalent).
\item
  Variables form one namespace (stage globals + sprite locals merged);
  every declared variable is initialized to \texttt{0}
  (Scratch\textquotesingle s default). Variable names are sanitized
  (\texttt{1st\ value} → \texttt{v\_1st\_value}) and kept collision-free
  (\texttt{print} → \texttt{print\_2}).
\end{itemize}

\subsubsection{5.2 Statements}\label{52-statements}

{\def\LTcaptype{none} % do not increment counter
\begin{longtable}[]{@{}ll@{}}
\toprule\noalign{}
Scratch & Python \\
\midrule\noalign{}
\endhead
\bottomrule\noalign{}
\endlastfoot
\texttt{set\ var\ to\ X} & \texttt{var\ =\ X} \\
\texttt{change\ var\ by\ X} & \texttt{var\ =\ \_wc\_add(var,\ X)} \\
\texttt{ask\ Q\ and\ wait} &
\texttt{\_\ =\ \_wc\_fmt(Q);\ \_answer\ =\ input()} (the question is
evaluated but never printed --- it is a stage bubble) \\
\texttt{say\ M} / \texttt{think\ M} & \texttt{print(\_wc\_fmt(M))} \\
\texttt{say\ M\ for\ N\ secs} & \texttt{print(\_wc\_fmt(M))} ---
duration ignored (stage timing), documented \\
\texttt{repeat\ N\ \{…\}} &
\texttt{for\ \_wc\_i\ in\ range(int(\_wc\_num(N))):} \\
\texttt{repeat\ until\ C\ \{…\}} & \texttt{while\ not\ C:} \\
\texttt{while\ C\ \{…\}} / \texttt{until\ C\ \{…\}} & \texttt{while\ C:}
/ \texttt{while\ not\ C:} \\
\texttt{forever\ \{…\}} & \texttt{while\ True:} (pair with
\texttt{stop\ all}) \\
\texttt{if\ C\ \{…\}} / \texttt{if\ C\ \{…\}\ else\ \{…\}} &
\texttt{if\ C:} / \texttt{if\ …\ else:} \\
\end{longtable}
}

\subsubsection{5.3 Expressions}\label{53-expressions}

Arithmetic is wrapped in \texttt{\_wc\_num} (Scratch treats non-numeric
values as 0 in arithmetic); comparisons use
\texttt{\_wc\_lt/\_wc\_gt/\_wc\_eq} (numbers compare numerically,
everything else case-insensitively --- Scratch\textquotesingle s rules);
\texttt{and/or/not} map directly; \texttt{answer} reads
\texttt{\_answer}.

\subsubsection{5.4 Semantic deltas compensated in the generated
preamble}\label{54-semantic-deltas-compensated-in-the-generated-preamble}

{\def\LTcaptype{none} % do not increment counter
\begin{longtable}[]{@{}ll@{}}
\toprule\noalign{}
Scratch behavior & Helper \\
\midrule\noalign{}
\endhead
\bottomrule\noalign{}
\endlastfoot
\texttt{1/0\ =\ Infinity}, \texttt{0/0\ =\ NaN} & \texttt{\_wc\_div} \\
\texttt{x\ mod\ 0\ =\ NaN}; \texttt{mod} sign follows the divisor &
\texttt{\_wc\_mod} \\
\texttt{round(2.5)\ =\ 3}, \texttt{round(-2.5)\ =\ -3} (halves away from
zero) & \texttt{\_wc\_round} (Python\textquotesingle s round is
banker\textquotesingle s) \\
\texttt{sin/cos/tan} in \textbf{degrees} &
\texttt{\_wc\_sin/\_wc\_cos/\_wc\_tan} \\
\texttt{sqrt(-1)\ =\ NaN} & \texttt{\_wc\_sqrt} \\
\texttt{say\ 3.0} prints \texttt{3}; \texttt{say\ true} prints
\texttt{true} & \texttt{\_wc\_fmt} (also \texttt{Infinity}/\texttt{NaN}
by name) \\
\end{longtable}
}

\begin{center}\rule{0.5\linewidth}{0.5pt}\end{center}

\subsection{6. Teacher guide}\label{6-teacher-guide}

Block problems slot into classes/teme like any problem. Suggested ladder
for the gimnaziu cohort:

\begin{enumerate}
\def\labelenumi{\arabic{enumi}.}
\tightlist
\item
  \textbf{Scratch, VERY\_EASY}: echo problems --- ask/say only (the ask
  spec\textquotesingle s \texttt{sensing\_askandwait} +
  \texttt{sensing\_answer}).
\item
  \textbf{Scratch}: sum of two numbers (variables +
  \texttt{operator\_add}), last digit (\texttt{operator\_mod}), counting
  loops (\texttt{control\_repeat}).
\item
  \textbf{Blockly (Python generator)}: same problems re-solved with
  blocks whose generated code they can \emph{read} --- the bridge
  lesson.
\item
  \textbf{Blockly → typed}: the learner re-submits the same problem in
  typed Python; the generated-code preview makes the mapping visible.
\end{enumerate}

The curated block sets are the pedagogical control: a digit-extraction
problem can allow only math/variables blocks, withholding string/text
blocks exactly like the wave-C \texttt{py-no-string-ops} feature ban
does for typed Python.

\begin{center}\rule{0.5\linewidth}{0.5pt}\end{center}

\subsection{7. Complexity hints}\label{7-complexity-hints}

\texttt{src/components/blocks/complexity.ts} reads the
\textbf{generated} code and surfaces advisory chips under the Blockly
editor:

\begin{itemize}
\tightlist
\item
  \texttt{nested-loops} --- a loop directly inside a loop: "if each loop
  runs n times, the body runs n² times --- for n = 10⁵ that is 10¹⁰
  steps and will time out".
\item
  \texttt{no-exit-forever} --- a \texttt{while\ True} with no
  \texttt{break}/\texttt{exit} in its body.
\end{itemize}

It is indentation-based, works for both generators\textquotesingle{}
output, and is a \textbf{hint, never a verdict} --- the judge still
measures real time. (Pure function; unit-tested in
\texttt{complexity.test.ts}.)

\begin{center}\rule{0.5\linewidth}{0.5pt}\end{center}

\subsection{8. Testing strategy}\label{8-testing-strategy}

{\def\LTcaptype{none} % do not increment counter
\begin{longtable}[]{@{}lll@{}}
\toprule\noalign{}
Layer & What & Where \\
\midrule\noalign{}
\endhead
\bottomrule\noalign{}
\endlastfoot
Transpiler & full mapping + semantic helpers, \textbf{generated code is
executed and its stdout asserted} &
\texttt{backend/problems/test\_scratch\_transpiler.py} (93 tests) \\
Submit gate & 400 paths: not a scratch problem, disallowed block,
unmapped block, bad JSON &
\texttt{backend/submissions/test\_scratch\_judge.py} \\
Judge hook & transpiled code compiled under \texttt{python3}; transpile
failure ⇒ COMPILE\_ERROR; effective limits (python3 multipliers;
\texttt{scratch} override wins) &
\texttt{backend/submissions/test\_scratch\_judge.py} \\
Authoring & \texttt{block\_config} engine validation, opcode allowlist
validation, \texttt{effective\_limits} scratch row &
\texttt{backend/submissions/test\_scratch\_judge.py} \\
Frontend units & toolbox construction/curation, variable extraction,
default project, complexity analyzer &
\texttt{src/components/blocks/*.test.ts} \\
E2E (mock) & playground Scratch tab; scratch problem renders editor;
submit POST payload is the project JSON under
\texttt{language:\ \textquotesingle{}scratch\textquotesingle{}}; a 400
surfaces its message & \texttt{e2e/scratch.spec.ts} (+ Blockly:
\texttt{e2e/blocks.spec.ts}) \\
E2E (live) & a real Scratch submission judged AC and a planted-WA judged
WA --- see the wave\textquotesingle s live acceptance record & live
stack, per-wave \\
\end{longtable}
}

The i18n parity sweep (\texttt{src/i18n.parity.test.ts}) guarantees
every new UI string exists in all ten locales --- the
engines\textquotesingle{} RO/EN labels are first-class, not
afterthoughts.

\begin{center}\rule{0.5\linewidth}{0.5pt}\end{center}

\subsection{9. Security notes}\label{9-security-notes}

\begin{itemize}
\tightlist
\item
  Scratch submissions are \textbf{data, not code}, until our own
  deterministic transpiler turns them into Python. No learner-supplied
  text is ever \texttt{eval}\textquotesingle d; the transpiler emits a
  fixed AST-shaped subset (identifiers are sanitized and
  collision-checked against builtins).
\item
  The three-gate defense (submit 400 → judge re-validation → registry
  execution) mirrors waves C/E: a hostile client skipping the API still
  cannot judge a project the API would reject.
\item
  The Scratch editor loads third-party libraries (LLK\textquotesingle s,
  BSD-3/GPL-class Scratch licenses --- the packages\textquotesingle{}
  \texttt{TRADEMARK} files govern the Scratch name/cat usage: we use
  them for the educational compatibility the trademark policy explicitly
  permits, and we do not imply endorsement). Default assets are bundled
  locally; a creator-uploaded project\textquotesingle s other costumes
  load from the official asset CDN \textbf{client-side only}.
\item
  Judge side is unchanged: scratch programs run inside the same
  hardened, network-isolated judge-runner as everything else.
\end{itemize}

Known v1 limitations (deliberate, documented): no variable-creation
modal (creator pre-declares variables in the starter); no custom blocks
(procedures), lists, or motion/pen semantics in Scratch problems (the
opcode allowlist is the contract --- unmapped opcodes are rejected with
a clear message); \texttt{say\ for\ seconds} ignores its duration.

\begin{center}\rule{0.5\linewidth}{0.5pt}\end{center}

\subsection{10. Troubleshooting}\label{10-troubleshooting}

{\def\LTcaptype{none} % do not increment counter
\begin{longtable}[]{@{}ll@{}}
\toprule\noalign{}
Symptom & Cause / fix \\
\midrule\noalign{}
\endhead
\bottomrule\noalign{}
\endlastfoot
Scratch editor shows "failed to load" & The lazy chunk or a bundled
asset failed --- check the network tab for
\texttt{/scratch-assets/*.svg} and \texttt{/scratch-blocks-media/}; a
blank stage alone is fine (no WebGL). \\
Stage blank but blocks work & WebGL unavailable; renderer init is
skipped by design (§2.1 step 2). Judging is unaffected. \\
400 "blocks with no judged mapping" & The project uses an opcode outside
§3.2\textquotesingle s table (e.g. motion or join) --- remove it or
extend the transpiler (next section). \\
400 "uses blocks this problem does not allow" & The opcode is mapped but
not in this problem\textquotesingle s \texttt{scratch\_blocks}. \\
COMPILE\_ERROR "could not be translated" & The judge\textquotesingle s
re-validation failed on the exact bytes --- same reasons as the 400s
(defense in depth). \\
Cat costume missing & Asset md5 not under
\texttt{public/scratch-assets/} and CDN unreachable; the bundled default
never needs the CDN. \\
Blockly "Show generated code" empty & No blocks run at top level ---
every Blockly program starts from placed blocks. \\
\end{longtable}
}

\begin{center}\rule{0.5\linewidth}{0.5pt}\end{center}

\subsection{11. Worked examples}\label{11-worked-examples}

\subsubsection{11.1 Blockly --- "Last digit" (generator
python3)}\label{111-blockly--last-digit-generator-python3}

Default starter (\texttt{variables\_set\ n\ ←\ wc\_input\_number} →
\texttt{print(math\_modulo(n,10))}) generates:

\begin{Shaded}
\begin{Highlighting}[]
\NormalTok{n }\OperatorTok{=} \BuiltInTok{int}\NormalTok{(}\BuiltInTok{input}\NormalTok{())}
\BuiltInTok{print}\NormalTok{(n }\OperatorTok{\%} \DecValTok{10}\NormalTok{)}
\end{Highlighting}
\end{Shaded}

Submitted as \texttt{language:\ python3} --- judged normally (testcase
\texttt{42} → \texttt{2}).

\subsubsection{11.2 Scratch --- "Echo"}\label{112-scratch--echo}

Problem: \texttt{engine:\ scratch}, allowlist
\texttt{event\_whenflagclicked,\ sensing\_askandwait,\ sensing\_answer,\ looks\_say};
testcase input \texttt{hello}, expected output \texttt{hello}. The
learner\textquotesingle s project (one flag script):

\begin{Shaded}
\begin{Highlighting}[]
\FunctionTok{\{}
  \DataTypeTok{"targets"}\FunctionTok{:} \OtherTok{[}
    \FunctionTok{\{} \DataTypeTok{"isStage"}\FunctionTok{:} \KeywordTok{true}\FunctionTok{,} \DataTypeTok{"name"}\FunctionTok{:} \StringTok{"Stage"}\FunctionTok{,} \DataTypeTok{"variables"}\FunctionTok{:} \FunctionTok{\{\},} \DataTypeTok{"blocks"}\FunctionTok{:} \FunctionTok{\{\},}
      \DataTypeTok{"costumes"}\FunctionTok{:} \OtherTok{[}\FunctionTok{\{} \DataTypeTok{"name"}\FunctionTok{:} \StringTok{"backdrop1"}\FunctionTok{,} \DataTypeTok{"md5ext"}\FunctionTok{:} \StringTok{"cd21514d0531fdffb22204e0ec5ed84a.svg"}\FunctionTok{,} \DataTypeTok{"dataFormat"}\FunctionTok{:} \StringTok{"svg"} \FunctionTok{\}}\OtherTok{]} \FunctionTok{\}}\OtherTok{,}
    \FunctionTok{\{} \DataTypeTok{"isStage"}\FunctionTok{:} \KeywordTok{false}\FunctionTok{,} \DataTypeTok{"name"}\FunctionTok{:} \StringTok{"Sprite1"}\FunctionTok{,} \DataTypeTok{"variables"}\FunctionTok{:} \FunctionTok{\{\},} \DataTypeTok{"blocks"}\FunctionTok{:} \FunctionTok{\{}
        \DataTypeTok{"hat"}\FunctionTok{:}  \FunctionTok{\{} \DataTypeTok{"opcode"}\FunctionTok{:} \StringTok{"event\_whenflagclicked"}\FunctionTok{,} \DataTypeTok{"next"}\FunctionTok{:} \StringTok{"ask"}\FunctionTok{,} \DataTypeTok{"topLevel"}\FunctionTok{:} \KeywordTok{true}\FunctionTok{,} \DataTypeTok{"fields"}\FunctionTok{:} \FunctionTok{\{\},} \DataTypeTok{"inputs"}\FunctionTok{:} \FunctionTok{\{\},} \DataTypeTok{"x"}\FunctionTok{:} \DecValTok{60}\FunctionTok{,} \DataTypeTok{"y"}\FunctionTok{:} \DecValTok{60} \FunctionTok{\},}
        \DataTypeTok{"ask"}\FunctionTok{:}  \FunctionTok{\{} \DataTypeTok{"opcode"}\FunctionTok{:} \StringTok{"sensing\_askandwait"}\FunctionTok{,} \DataTypeTok{"next"}\FunctionTok{:} \StringTok{"say"}\FunctionTok{,} \DataTypeTok{"parent"}\FunctionTok{:} \StringTok{"hat"}\FunctionTok{,}
                  \DataTypeTok{"inputs"}\FunctionTok{:} \FunctionTok{\{} \DataTypeTok{"QUESTION"}\FunctionTok{:} \OtherTok{[}\DecValTok{1}\OtherTok{,} \OtherTok{[}\DecValTok{10}\OtherTok{,} \StringTok{"line?"}\OtherTok{]} \OtherTok{]} \FunctionTok{\},} \DataTypeTok{"fields"}\FunctionTok{:} \FunctionTok{\{\}} \FunctionTok{\},}
        \DataTypeTok{"say"}\FunctionTok{:}  \FunctionTok{\{} \DataTypeTok{"opcode"}\FunctionTok{:} \StringTok{"looks\_say"}\FunctionTok{,} \DataTypeTok{"parent"}\FunctionTok{:} \StringTok{"ask"}\FunctionTok{,} \DataTypeTok{"next"}\FunctionTok{:} \KeywordTok{null}\FunctionTok{,}
                  \DataTypeTok{"inputs"}\FunctionTok{:} \FunctionTok{\{} \DataTypeTok{"MESSAGE"}\FunctionTok{:} \OtherTok{[}\DecValTok{1}\OtherTok{,} \StringTok{"ans"}\OtherTok{]} \FunctionTok{\},} \DataTypeTok{"fields"}\FunctionTok{:} \FunctionTok{\{\}} \FunctionTok{\},}
        \DataTypeTok{"ans"}\FunctionTok{:}  \FunctionTok{\{} \DataTypeTok{"opcode"}\FunctionTok{:} \StringTok{"sensing\_answer"}\FunctionTok{,} \DataTypeTok{"fields"}\FunctionTok{:} \FunctionTok{\{\}} \FunctionTok{\}}
      \FunctionTok{\},}
      \DataTypeTok{"costumes"}\FunctionTok{:} \OtherTok{[}\FunctionTok{\{} \DataTypeTok{"name"}\FunctionTok{:} \StringTok{"cat1"}\FunctionTok{,} \DataTypeTok{"md5ext"}\FunctionTok{:} \StringTok{"bcf454acf82e4504149f7ffe07081dbc.svg"}\FunctionTok{,} \DataTypeTok{"dataFormat"}\FunctionTok{:} \StringTok{"svg"} \FunctionTok{\}}\OtherTok{]} \FunctionTok{\}}
  \OtherTok{]}
\FunctionTok{\}}
\end{Highlighting}
\end{Shaded}

Submitted as \texttt{language:\ scratch}; the judge transpiles
(abridged):

\begin{Shaded}
\begin{Highlighting}[]
\NormalTok{\_ }\OperatorTok{=}\NormalTok{ \_wc\_fmt(}\StringTok{"line?"}\NormalTok{)}
\NormalTok{\_answer }\OperatorTok{=} \BuiltInTok{input}\NormalTok{()}
\BuiltInTok{print}\NormalTok{(\_wc\_fmt(\_answer))}
\end{Highlighting}
\end{Shaded}

Input \texttt{hello} → output \texttt{hello} → \textbf{AC}. (Planted-WA
acceptance: a project that says \texttt{"goodbye"} instead →
\textbf{WA}.)

\subsubsection{11.3 Scratch --- "Sum of two
numbers"}\label{113-scratch--sum-of-two-numbers}

Two \texttt{ask} blocks, \texttt{set\ total\ to\ (answer)\ +\ (answer)},
\texttt{say\ total}:

\begin{Shaded}
\begin{Highlighting}[]
\NormalTok{\_ }\OperatorTok{=}\NormalTok{ \_wc\_fmt(}\StringTok{"a?"}\NormalTok{)}
\NormalTok{\_answer }\OperatorTok{=} \BuiltInTok{input}\NormalTok{()}
\NormalTok{total }\OperatorTok{=}\NormalTok{ \_wc\_num(\_answer)}
\NormalTok{\_ }\OperatorTok{=}\NormalTok{ \_wc\_fmt(}\StringTok{"b?"}\NormalTok{)}
\NormalTok{\_answer }\OperatorTok{=} \BuiltInTok{input}\NormalTok{()}
\NormalTok{total }\OperatorTok{=}\NormalTok{ (\_wc\_num(total) }\OperatorTok{+}\NormalTok{ \_wc\_num(\_answer))}
\BuiltInTok{print}\NormalTok{(\_wc\_fmt(total))}
\end{Highlighting}
\end{Shaded}

\texttt{3} + \texttt{4} prints \texttt{7} --- and \texttt{say\ (7/2)}
would print \texttt{3.5}, while \texttt{say\ (7/2*2)} prints \texttt{7},
matching Scratch\textquotesingle s rendering exactly.

\begin{center}\rule{0.5\linewidth}{0.5pt}\end{center}

\subsection{12. Screenshots (described)}\label{12-screenshots-described}

\begin{enumerate}
\def\labelenumi{\arabic{enumi}.}
\tightlist
\item
  \textbf{The /blocks playground} --- engine switcher (Blockly
  \textbar{} Scratch) top left; Blockly workspace with the beginner
  toolbox (categories: Logic, Loops, Math, Text, Input, Lists,
  Variables, Functions); the generated-Python preview pane below; the
  "Run on the judge" panel with a stdin box.
\item
  \textbf{The playground, Scratch engine} --- split view: scratch-blocks
  workspace (Scratch-shaped palette: Events, Control, Looks, Sensing,
  Operators, Variables) on the left; on the right the green-flag/stop
  toolbar, the white 4:3 stage with the orange Scratch cat, and the
  console strip that just printed \texttt{Hello!}.
\item
  \textbf{A scratch problem page} --- the statement card; the submit
  card titled "Solve with Scratch" with the same split editor; the
  effective-limits line showing \texttt{scratch:\ 3000\ ms\ /\ 384\ MB}
  (python3 multipliers); the verdict panel after judging.
\item
  \textbf{Complexity hints} --- a Blockly workspace with two nested
  repeat blocks; under the editor a ⚠ chip: "Nested loops: \ldots{} n²
  \ldots{} will time out."
\item
  \textbf{Creator form} --- "Block programming (Blockly / Scratch)"
  checked; the engine dropdown offering both; the Scratch arm showing
  the starter-project JSON box and the allowed-blocks opcode box; the
  languages grid with "Scratch (blocks)" ticked showing
  \texttt{(3000\ ms\ /\ 384\ MB)}.
\end{enumerate}

\begin{center}\rule{0.5\linewidth}{0.5pt}\end{center}

\subsection{13. Internationalization}\label{13-internationalization}

\begin{itemize}
\tightlist
\item
  \textbf{UI chrome} (\texttt{Solve\ with\ Scratch}, category names,
  hints, buttons): the site\textquotesingle s ten-locale catalogs
  (\texttt{src/locales/*.json}, keys \texttt{sbx\_*},
  \texttt{blocks\_*}, \texttt{problem\_form\_*}); the parity test pins
  completeness.
\item
  \textbf{Blockly block labels}: official Blockly message packs
  (\texttt{blockly/msg/ro}, \texttt{blockly/msg/en}) + the two
  site-block messages (\texttt{WC\_BLOCK\_MESSAGES\_RO/EN} in
  \texttt{BlockEditor.tsx}).
\item
  \textbf{Scratch block labels}: official scratch-l10n
  \texttt{editor/blocks/ro.json} applied to \texttt{Blockly.Msg} (the
  scratch-gui mechanism); EN is baked into scratch-blocks. The category
  headers localize through the site catalog.
\end{itemize}

\begin{center}\rule{0.5\linewidth}{0.5pt}\end{center}

\subsection{14. Extending the system}\label{14-extending-the-system}

\textbf{Add a Blockly category/block}: add the block types to
\texttt{DEFAULT\_TOOLBOX\_XML} (or a problem\textquotesingle s toolbox
XML) --- the official generators already know the standard blocks. For
new \emph{site} blocks, follow \texttt{wcBlocks.ts}: JSON definition +
generator snippets for BOTH Python and JavaScript, register them in
\texttt{registerBlockTypes}, and test the generated code contract
headlessly (see \texttt{BlockEditor.test.ts}).

\textbf{Add a Scratch opcode} (e.g. \texttt{operator\_join}): extend
\texttt{STATEMENT\_OPCODES}/\texttt{REPORTER\_OPCODES} + the emitter
branch in \texttt{scratch\_transpiler.py}, add the opcode to
\texttt{OPCODE\_CATEGORIES} in \texttt{ScratchEditor.tsx} (palette), and
cover the mapping with a test that \textbf{executes} the generated code
(\texttt{test\_scratch\_transpiler.py} is the pattern). Keep the
contract: opcodes without an honest mapping are rejected, never
approximated.

\textbf{Upgrade the Scratch stack}: bump the pinned versions (they are
pinned to the scratch-gui-5.3-compatible set), re-run
\texttt{cp\ -r\ node\_modules/scratch-blocks/media\ public/scratch-blocks-media/},
and re-verify the wiring glue (§2.1) --- scratch-vm\textquotesingle s
\texttt{blockListener} / \texttt{emitWorkspaceUpdate} / \texttt{toJSON}
API is the surface we depend on.

\end{document}
